\section{PLANTEAMIENTO DEL PROBLEMA}


Se refiere al interrogante que lleva al investigador a buscar respuestas concretas. Es la definición del problema que aborda con la investigación. 


\subsection{Antecedentes}
Los antecedentes son las investigaciones que se han realizado previamente y que guardan una relación histórica con el tema de investigación actual. 


\newpage

% ==================================================

\section{JUSTIFICACIÓN}


Responde a los interrogantes del por qué se desea conocer el tema y por qué se seleccionó, así como cuál es el aporte que tendrá el texto a la ciencia. 

No abuse del uso de \textit{cursivas} o \textbf{negritas} dentro del texto, úselas muy moderadamente, por lo general saturan y dificultan la lectura del documento. Utilice cursivas en casos muy particulares como géneros y especies (\textit{Tyrannus melancholicus}), términos químicos (\textit{kr}), letras griegas ($\beta$) y algunos títulos y subtítulos. El entorno matemático de \LaTeX \ muestra las variables y símbolos en letra cursiva, asegurese de hacer lo mismo al referenciarla en el documento, Utilice \textbf{negritas} en algunos títulos de capítulos y subcapítulos, algunos datos de tablas o enfatizar aspectos muy particulares. El uso de \underline{texto subrayado} no se recomienda en normas IEEE.

Utilice moderadamente el uso de abreviaturas, se prefiere que el texto sea más largo y claro que corto y confuso para el lector. Por ejemplo, IEEE puede significar Institute of Electrical and Electronics Engineers o Instituto Español de Estudios Estratégicos. Sin embargo, las abreviaturas pueden ser útiles en casos como la repetición continua en un mismo párrafo.

Prefiera las comillas “inglesas” y ‘sencillas’ por sobre las «latinas» o «españolas».

\begin{description}
    \item[Características:] texto descriptivo.
    \item[Propiedades:] texto descriptivo.
    \item[Estructura:] texto descriptivo.
\end{description}